\chapter{باقیمانده‌های مرکب درجۀ \en{n}}
\label{app:dcr}
\section{تعریف مسئله}
اگر $n=pq$ حاصل‌ضرب دو عدد اول باشد، یک عضو $z\in\Z_{n^2}^{*}$ باقیماندۀ درجۀ $n$ است هرگاه عضوی $u$ وجود داشته باشد که
\begin{equation}
z=u^n\pmod{n^2}.
\end{equation}
مسئلۀ تصمیمی باقیماندۀ مرکب، تشخیص عضویت یک عنصر در این زیرگروه است. دشواری این مسئله، از مبانی تحلیل رمزگذاری پاییه است \cite{ref179}. دانستن عوامل $n$ ابزارهای محاسبۀ ساختار گروه و حل مسئله را فراهم می‌کند؛ با این حال، نباید بدون اثبات، دشواری این فرض را معادل دقیق دشواری تجزیه دانست.

\section{رابطه با رمزگذاری پاییه}
در متن رمز $c=g^m u^n\pmod{n^2}$، مؤلفۀ $g^m$ پیام را در کلاس باقیمانده قرار می‌دهد و $u^n$ تصادفی‌سازی در زیرگروه مربوط ایجاد می‌کند. با اعمال توان خصوصی $\lambda$، مؤلفۀ تصادفی حذف می‌شود. برای $g=1+n$، بسط دوجمله‌ای در پیمانۀ $n^2$ می‌دهد:
\begin{equation}
(1+n)^a\equiv1+an\pmod{n^2},\qquad
L((1+n)^a\bmod n^2)=a\pmod n.
\end{equation}
بنابراین اگر شرایط وارون‌پذیری برقرار باشند، اعمال $L$ و ضرب در وارون مناسب، پیام را بازمی‌گرداند. در فضای بزرگ‌تر دامگارد\lr{--}یوریک، جمله‌های مرتبۀ دوم و بالاتر باقی می‌مانند؛ همین نکته علت نادرست بودن تعمیم بی‌واسطۀ فرمول سادۀ پاییه برای تمام $s$ است.

\section{حدود فرض و حملات مرتبط}
طرح جمعی ذاتاً اجازۀ تغییر متن رمز با یک عملگر جبری را می‌دهد. این ویژگی برای محاسبه لازم است، اما به معنای مقاومت خودکار در برابر حملۀ فعال یا متن رمز انتخاب‌شده نیست. احراز اصالت محموله، سیاست پذیرش، جلوگیری از بازپخش و محدود کردن دسترسی به رمزگشا باید در مدل کاربردی بررسی شوند.

استفاده از یک تبدیل عمومی برای دستیابی به امنیت قوی‌تر نیز لزوماً همریختی را حفظ نمی‌کند؛ در بسیاری از تبدیل‌ها، تغییر جبری متن رمز باعث نامعتبر شدن لایۀ اضافه می‌شود. ازاین‌رو طراحی امنِ یک سامانۀ محاسباتی نیازمند سازگاری لایۀ احراز اصالت با تابع مورد نیاز است. در ساختار فصل چهارم، امضا پیش از تجمیع بررسی می‌شود و مؤلفۀ امضا در جبر پایۀ رمز دخالت داده نمی‌شود.

\section{پیاده‌سازی و آیندۀ پارامترها}
تولید تصادف امن، کنترل نسبتاً اول بودن، بررسی عضو معتبر و نگهداری محرمانۀ عوامل پیمانه، بخش‌های ضروری پیاده‌سازی‌اند. نشت کانال جانبی یا افشای یکی از عوامل می‌تواند فرض ریاضیِ مدل ایده‌آل را بی‌اثر کند. انتخاب اعداد اول امن ممکن است برای برخی ساختارهای گروهی مفید باشد، اما جایگزین کامل شرایط معتبر پارامترهای هر طرح نیست.

توانایی کوانتومیِ کافی برای تجزیۀ $n$، طرح‌های متکی بر این خانوادۀ ساختارها را تهدید می‌کند. اندازه‌های آزمایشی این پایان‌نامه برای نتیجه‌گیری دربارۀ مقاومت بلندمدت مناسب نیستند. پژوهش دربارۀ ساختارهای پساکوانتومی و ارزیابی هزینۀ محاسباتی آن‌ها یک مسیر مستقل برای توسعۀ سامانه‌های حساس است.
